\documentclass[12pt]{article}

\usepackage[margin=1in]{geometry}
\usepackage{amsmath,amsthm,amssymb,mathtools}
\usepackage{booktabs}
\usepackage{array}
\usepackage{microtype}
\usepackage{parskip}
\usepackage{fancyhdr}
\usepackage{titlesec}
\usepackage[T1]{fontenc}
\usepackage{lmodern}
\usepackage[utf8]{inputenc}
\usepackage[hidelinks]{hyperref}

% -----------------------------------------------------------
% Header / footer
% -----------------------------------------------------------
\pagestyle{fancy}
\fancyhf{}
\lhead{\small N-Series Acceleration for Polygonal $\pi$}
\rhead{\small Wayne Baker}
\cfoot{\thepage}
\setlength{\headheight}{14.5pt}

% -----------------------------------------------------------
% Section formatting
% -----------------------------------------------------------
\titleformat{\section}
  {\large\bfseries}{\thesection.}{0.5em}{}
\titleformat{\subsection}
  {\normalsize\bfseries}{\thesubsection.}{0.5em}{}

% -----------------------------------------------------------
% Theorem environments
% -----------------------------------------------------------
\newtheorem{theorem}{Theorem}[section]
\newtheorem{proposition}[theorem]{Proposition}
\newtheorem{lemma}[theorem]{Lemma}
\newtheorem{corollary}[theorem]{Corollary}
\theoremstyle{definition}
\newtheorem{definition}[theorem]{Definition}
\theoremstyle{remark}
\newtheorem{remark}[theorem]{Remark}

% -----------------------------------------------------------
% Commands
% -----------------------------------------------------------
\newcommand{\Rop}{\mathcal{R}}
\newcommand{\Rmulti}{\mathcal{E}}
\newcommand{\abs}[1]{\left|#1\right|}
\newcommand{\Order}{\mathcal{O}}

% -----------------------------------------------------------
\begin{document}

\begin{center}
  {\LARGE\bfseries Certified Constructive N-Series Acceleration}\\[4pt]
  {\LARGE\bfseries for Polygonal Approximation of $\pi$}\\[16pt]
  {\large Wayne Baker}\\[3pt]
  {\normalsize\itshape Independent Researcher, Newfoundland and Labrador, Canada}\\[3pt]
  {\normalsize November 15, 2026}
\end{center}

\begin{abstract}
We study a tripling-based geometric refinement and scale-cancellation framework
for Archimedean polygonal approximation of $\pi$.  For the normalized
semiperimeter
\[
  P_N=N\sin\!\left(\frac{\pi}{N}\right)
\]
of the regular inscribed $N$-gon in the unit circle, the classical even-power
expansion implies an error of order $N^{-2}$.  Consequently, raw tripling
$N\mapsto 3N$ reduces the leading error asymptotically by a factor of $9$,
whereas raw doubling reduces it by a factor of $4$.  After $k$ matched
refinement stages, the asymptotic same-stage advantage is $(9/4)^k$; at
$k=12$ the exact polygonal reference values differ by a factor of approximately
$1.6834\times10^4$.

The principal N-Series acceleration is not tripling alone.  It is the
scale-cancellation operator
\[
  \Rop_m[P]_N=\frac{m^2P_{mN}-P_N}{m^2-1},
\]
which cancels the leading $N^{-2}$ term and yields fourth-order error.  For
$m=3$, the asymptotic error reduction under one further tripling is therefore
$3^4=81$, compared with $2^4=16$ for a fourth-order doubling sequence.  More
general multiscale combinations cancel successive even-power error modes by
standard Richardson--Romberg moment conditions.

A separate analysis distinguishes three levels that were conflated in earlier
accounts: the ideal regular-polygon sequence, a $\pi$-free algebraic chord
recurrence, and a straightedge-and-compass approximate-trisection lane.  If
$c_N=\sin(\pi/N)$, the triple-angle identity determines $c_{3N}$ as the
smallest positive root of $4u^3-3u+c_N=0$; for $0<c_N<1$ the cubic has two
positive roots, and the selected branch is the smaller.  This recurrence generates the polygonal data
without inserting $\pi$, but a numerical cubic solve is not an exact Euclidean
construction of the regular $18$-gon.  The constructible lane instead produces
a finite chord approximation whose angular residual must be carried as a
separate error term.  The resulting geometric limit, denoted $\pi_g$, equals
$\pi$ because the underlying polygonal sequence converges to $\pi$ and the
N-Series transformations preserve that limit.  Preserving the \emph{accelerated}
fourth-order law along the constructible lane is a stronger requirement: we show
that a fixed number of cubic refinements per tripling level does not suffice,
because the refinement acts only on the fresh local defect and not on the
inherited angular error, and we replace it with a target-dependent schedule
whose exact-rational per-level budget is fixed in advance of the ladder.

The finite canonical scale-$3$ computation is now accompanied by a rigorous,
oracle-free error certificate.  Eleven polygon inputs are enclosed by exact
rational cubic-root intervals, propagated through the complete Richardson
tableau with exact rational weights, and combined with an analytic alternating-
series remainder using only the rational inequality $\pi<22/7$.  For the
specified $120$-decimal working-precision run, the resulting bound is
\[
  \abs{\widehat\pi_g-\pi}
  <3.406285758314118\times10^{-81},
\]
which certifies at least $80$ decimal places of absolute accuracy.  A separate
million-digit comparison with an independently generated Chudnovsky expansion
remains a large-scale implementation test; it is not claimed as a million-digit
formal enclosure.
\end{abstract}

% -----------------------------------------------------------
\section{Introduction}

Archimedean approximation of $\pi$ replaces the circle by regular inscribed and
circumscribed polygons~\cite{archimedes}.  In the familiar doubling scheme, a regular hexagon is
refined through
\[
  6,12,24,48,\ldots,
\]
and the leading perimeter error decreases by a factor approaching $4$ at each
step.  A tripling schedule instead uses the scale sequence
\[
  6,18,54,162,\ldots,
\]
for which the corresponding leading-error factor approaches $9$.

Two mathematically distinct mechanisms must be separated.

First, \emph{raw scale refinement} replaces $N$ by $mN$.  If the leading error
is proportional to $N^{-2}$, this operation reduces that term by $m^2$.
Tripling is therefore stronger than doubling per refinement level, although it
also increases the side count more rapidly.

Second, \emph{scale cancellation} forms a weighted combination of approximants
at two or more resolutions.  The weights are chosen to annihilate the leading
terms in an even-power asymptotic expansion.  The two-scale operator used here
is a Richardson extrapolation with scale ratio $m$; the term \emph{Baker
N-Series} refers to its organization along geometric polygonal scale families,
especially the tripling family $N,3N,9N,\ldots$, together with a finite
constructible interpretation of the underlying chord stages.

A further distinction is essential.  The analytic reference value
$P_{18}=18\sin(\pi/18)$ belongs to an exact regular $18$-gon, which is not
classically constructible.  The algebraic computational lane obtains its chord
by selecting a root of a cubic, whereas Baker's geometric lane produces a
straightedge-and-compass approximation to the required central angle.  Thus a
publication-safe analysis must carry two errors:
\[
  \text{polygonal discretization error}
  \quad\text{and}\quad
  \text{constructive angular residual}.
\]
The first is governed by the even-power expansion; the second is governed by
the local cubic refinement law developed in the companion trisection papers.

The contributions of this paper are as follows:
\begin{enumerate}
  \item an exact comparison of raw doubling and tripling for the inscribed
        polygonal sequence;
  \item a fourth-order N-Series cancellation theorem and its multiscale
        extension;
  \item a $\pi$-free algebraic chord-tripling recurrence derived from the
        triple-angle identity;
  \item a perturbation estimate specifying how accurately the constructible
        central angle must approximate the ideal angle, together with an
        exact-rational refinement schedule that meets it;
  \item a proof that the geometric limit $\pi_g$ equals $\pi$;
  \item an exact-rational finite-value enclosure for the canonical eleven-input
        scale-$3$ Richardson tableau;
  \item numerical tables separating raw refinement, scale cancellation,
        certified finite-value accuracy, and arbitrary-precision implementation
        verification;
  \item a precise statement of what the million-digit computation does and
        does not establish.
\end{enumerate}

% -----------------------------------------------------------
\section{The exact polygonal reference sequence}

\subsection{Normalized semiperimeter}

For a regular inscribed $N$-gon in the unit circle, the side length is
$2\sin(\pi/N)$.  Its perimeter is therefore $2N\sin(\pi/N)$.  It is convenient
to normalize by a factor of $2$ and define
\begin{equation}
  P_N:=N\sin\!\left(\frac{\pi}{N}\right).
  \label{eq:PN}
\end{equation}
Then $P_N<\pi$ and $P_N\to\pi$ as $N\to\infty$.

Expanding the sine function at the origin gives
\begin{equation}
\begin{aligned}
  P_N
  &=\pi-\frac{\pi^3}{6N^2}
       +\frac{\pi^5}{120N^4}
       -\frac{\pi^7}{5040N^6}
       +\Order(N^{-8}).
\end{aligned}
\label{eq:PNexpansion}
\end{equation}
Hence the positive error
\begin{equation}
  E_N:=\pi-P_N
  \label{eq:EN}
\end{equation}
satisfies
\begin{equation}
  E_N=\frac{\pi^3}{6N^2}
      -\frac{\pi^5}{120N^4}
      +\frac{\pi^7}{5040N^6}
      +\Order(N^{-8}).
  \label{eq:ENexpansion}
\end{equation}

\subsection{Raw scale refinement}

\begin{theorem}[General raw scale law]
Let
\[
  Q_N=L+cN^{-p}+o(N^{-p}),
  \qquad c\ne0,\quad p>0.
\]
For every fixed scale factor $m>1$,
\begin{equation}
  \frac{Q_{mN}-L}{Q_N-L}\longrightarrow m^{-p}
  \qquad (N\to\infty).
  \label{eq:general-scale-law}
\end{equation}
Equivalently, the magnitude of the leading error is reduced asymptotically by
the factor $m^p$.
\end{theorem}

\begin{proof}
From the assumed expansion,
\[
  Q_{mN}-L=c(mN)^{-p}+o(N^{-p})
          =cm^{-p}N^{-p}+o(N^{-p}).
\]
Division by $Q_N-L=cN^{-p}+o(N^{-p})$ gives
\eqref{eq:general-scale-law}.
\end{proof}

Applying the theorem to \eqref{eq:ENexpansion} with $p=2$ gives the following.

\begin{corollary}[Doubling versus tripling]
For the inscribed-polygon error $E_N$,
\[
  \frac{E_{2N}}{E_N}\longrightarrow\frac14,
  \qquad
  \frac{E_{3N}}{E_N}\longrightarrow\frac19.
\]
Thus raw doubling reduces the leading error by a factor of $4$, whereas raw
tripling reduces it by a factor of $9$.
\end{corollary}

If both schedules begin from the same $N_0$ and are compared after $k$ matched
refinement stages, then
\begin{equation}
  \frac{E_{2^kN_0}}{E_{3^kN_0}}
  \sim\left(\frac{3}{2}\right)^{2k}
  =\left(\frac94\right)^k.
  \label{eq:same-stage-ratio}
\end{equation}
At $k=12$,
\[
  \left(\frac94\right)^{12}
  =16834.112196028232\ldots.
\]
This is a same-stage error comparison, not a wall-clock speedup claim.

\begin{remark}[Digits gained per raw scale level]
For an $N^{-2}$ error law, doubling gains asymptotically
\[
  \log_{10}(4)=0.6020599913279624\ldots
\]
decimal digits per level, whereas tripling gains
\[
  \log_{10}(9)=0.9542425094393249\ldots.
\]
The ratio is
\[
  \frac{\log 9}{\log 4}
  =\frac{\log 3}{\log 2}
  =1.584962500721156\ldots.
\]
Thus tripling gains about $58.50\%$ more digits per scale level and requires
about $36.91\%$ fewer levels to reach the same asymptotic error.  The cost per
level, however, need not be the same.
\end{remark}

% -----------------------------------------------------------
\section{N-Series scale cancellation}

\subsection{The two-scale operator}

The raw factor $9$ is not the principal superconvergence mechanism.  The
N-Series acceleration combines two scales so that the leading error term
vanishes.

\begin{definition}[Two-scale N-Series operator]
For a fixed $m>1$, define
\begin{equation}
  \Rop_m[Q]_N
  :=\frac{m^2Q_{mN}-Q_N}{m^2-1}.
  \label{eq:Rop}
\end{equation}
\end{definition}

\begin{theorem}[Cancellation of the $N^{-2}$ mode]
Suppose
\begin{equation}
  Q_N=L+c_2N^{-2}+c_4N^{-4}+c_6N^{-6}+\Order(N^{-8}).
  \label{eq:general-even-expansion}
\end{equation}
Then
\begin{equation}
  \Rop_m[Q]_N
  =L-\frac{c_4}{m^2}N^{-4}
    -\frac{m^2+1}{m^4}c_6N^{-6}
    +\Order(N^{-8}).
  \label{eq:Rop-expansion}
\end{equation}
In particular,
\[
  \Rop_m[Q]_N=L+\Order(N^{-4}).
\]
\end{theorem}

\begin{proof}
Substitute \eqref{eq:general-even-expansion} at the scales $N$ and $mN$ into
\eqref{eq:Rop}.  The constant term has coefficient
$(m^2-1)/(m^2-1)=1$, while the $N^{-2}$ coefficient is
\[
  \frac{m^2m^{-2}-1}{m^2-1}=0.
\]
The $N^{-4}$ coefficient is
\[
  \frac{m^2m^{-4}-1}{m^2-1}
  =\frac{m^{-2}-1}{m^2-1}
  =-\frac1{m^2},
\]
and the $N^{-6}$ coefficient is
\[
  \frac{m^2m^{-6}-1}{m^2-1}
  =\frac{m^{-4}-1}{m^2-1}
  =-\frac{m^2+1}{m^4}.
\]
This proves \eqref{eq:Rop-expansion}.
\end{proof}

For \eqref{eq:PNexpansion}, $c_4=\pi^5/120$.  Therefore
\begin{equation}
  \Rop_m[P]_N
  =\pi-\frac{\pi^5}{120m^2N^4}+\Order(N^{-6}).
  \label{eq:Rop-P}
\end{equation}
The tripling operator is
\begin{equation}
  \Rop_3[P]_N=\frac{9P_{3N}-P_N}{8}
             =\pi-\frac{\pi^5}{1080N^4}+\Order(N^{-6}).
  \label{eq:R3}
\end{equation}
Thus one further tripling reduces its leading fourth-order error by
$3^4=81$, whereas one further doubling in a fourth-order doubling sequence
reduces the corresponding leading error by $2^4=16$.  The per-stage
error-factor ratio is
\[
  \frac{81}{16}=5.0625,
\]
while the ratio of decimal digits gained per scale level remains
$\log 3/\log 2$.

\subsection{Higher N-Series columns}

The same idea extends to several scales.  The algebra is the standard
Richardson--Romberg cancellation mechanism~\cite{richardson-gaunt,romberg}; the
geometric content lies in the particular polygonal scale family and its finite
constructible realization.

\begin{proposition}[Multiscale even-power cancellation]
Suppose
\begin{equation}
  Q_N=L+\sum_{r=1}^{s}c_{2r}N^{-2r}
      +\Order(N^{-2s-2}).
  \label{eq:multi-expansion}
\end{equation}
Fix $m>1$.  There is a unique weight vector
$w_0,\ldots,w_s$ satisfying
\begin{equation}
  \sum_{j=0}^{s}w_j=1,
  \qquad
  \sum_{j=0}^{s}w_jm^{-2rj}=0
  \quad (r=1,\ldots,s).
  \label{eq:moment-conditions}
\end{equation}
For these weights,
\begin{equation}
  \Rmulti_{m,s}[Q]_N
  :=\sum_{j=0}^{s}w_jQ_{m^jN}
  =L+\Order(N^{-2s-2}).
  \label{eq:multi-operator}
\end{equation}
\end{proposition}

\begin{proof}
The linear system \eqref{eq:moment-conditions} is a Vandermonde system in the
distinct nodes
\[
  1,\ m^{-2},\ m^{-4},\ \ldots,\ m^{-2s},
\]
and therefore has a unique solution.  Substitution of
\eqref{eq:multi-expansion} into \eqref{eq:multi-operator} shows that the
normalization preserves $L$ and that the moment conditions cancel the terms
$N^{-2},\ldots,N^{-2s}$.
\end{proof}

\begin{remark}[Novelty and attribution]
\label{rem:novelty}
The cancellation identities, the even-power expansion of $N\sin(\pi/N)$, and
Richardson extrapolation applied to polygonal perimeters are all classical; the
last is already present in Richardson and Gaunt~\cite{richardson-gaunt}, and the
even-power expansion is considerably older.  No priority is claimed for any of
them.  The contributions claimed here are instead the $\pi$-free cubic chord
ladder of Section~\ref{subsec:chord-recurrence} with its exact branch
identification (Theorem~\ref{thm:triple-angle-identification}); the oracle-free
exact-rational finite-value certificate of
Theorem~\ref{thm:certified-canonical-value}; the perturbation calculus of
Lemma~\ref{lem:transfer} together with the inheritance identity and
target-dependent budget of Section~\ref{subsec:ladder}, which separate the
constructive angular residual from the scale-truncation error and fix the
refinement depth required to preserve each asymptotic law; and the companion
Lean~4 formalization of Section~\ref{sec:lean}.
\end{remark}

% -----------------------------------------------------------
\section{Algebraic and constructible tripling lanes}

\subsection{A $\pi$-free algebraic chord recurrence}
\label{subsec:chord-recurrence}

Set
\begin{equation}
  c_N:=\sin\!\left(\frac{\pi}{N}\right),
  \qquad P_N=Nc_N.
  \label{eq:chord-half-length}
\end{equation}
The identity
\[
  \sin(3x)=3\sin x-4\sin^3x
\]
with $x=\pi/(3N)$ gives
\begin{equation}
  4u^3-3u+c_N=0,
  \qquad u=c_{3N}.
  \label{eq:cubic-tripling}
\end{equation}
Starting from an algebraic seed such as
\[
  c_6=\frac12
  \qquad\text{or}\qquad
  c_3=\frac{\sqrt3}{2},
\]
one may generate the entire tripling ladder by repeatedly selecting the
smallest positive root of \eqref{eq:cubic-tripling}.  No decimal expansion of $\pi$ is
required to generate the chord data; $\pi$ enters only afterward as an
external reference for measuring the error.

\begin{remark}[Algebraic does not mean Euclidean-constructible]
Equation~\eqref{eq:cubic-tripling} supplies an exact algebraic relation, but a
general cubic root is not obtainable by straightedge and compass.  The
$\pi$-free production script is therefore an algebraic numerical realization
of exact chord tripling, not a claim that every regular $3N$-gon is classically
constructible.  The approximate-trisection lane below is the genuinely finite
straightedge-and-compass realization.
\end{remark}

\begin{proposition}[Unique selected cubic branch]
\label{prop:unique-cubic-branch}
For every $0<c\le\tfrac12$, the polynomial
\[
  f_c(u)=4u^3-3u+c
\]
has exactly one zero in the interval $(0,c/2)$.
\end{proposition}

\begin{proof}
At the endpoints,
\[
  f_c(0)=c>0,
  \qquad
  f_c\!\left(\frac c2\right)=\frac{c(c^2-1)}2<0.
\]
Moreover,
\[
  f_c'(u)=12u^2-3<0
\]
for $0\le u\le c/2\le1/4$.  Thus $f_c$ is strictly decreasing on a
sign-changing interval and has exactly one zero there.
\end{proof}

\begin{theorem}[Exact identification of the algebraic tripling ladder]
\label{thm:triple-angle-identification}
Let $c_0=1/2$, and define $c_{L+1}$ to be the unique zero of
\[
  4u^3-3u+c_L=0
\]
in $(0,c_L/2)$.  Then, for every $L\ge0$,
\[
  c_L=\sin\!\left(\frac{\pi}{6\cdot3^L}\right).
\]
Consequently, with $N_L=6\cdot3^L$,
\[
  N_Lc_L=N_L\sin\!\left(\frac{\pi}{N_L}\right)=P_{N_L}.
\]
\end{theorem}

\begin{proof}
The base case is $c_0=1/2=\sin(\pi/6)$.  Assume
$c_L=\sin x$ with $x=\pi/(6\cdot3^L)$, and set $u=\sin(x/3)$.  The triple-angle
identity gives $4u^3-3u+c_L=0$.  Also,
\[
  c_L-2u=(3u-4u^3)-2u=u(1-4u^2)>0,
\]
so $u\in(0,c_L/2)$.  Proposition~\ref{prop:unique-cubic-branch} makes $u$ the
unique selected root, and induction completes the proof.
\end{proof}

\subsection{Approximate trisection}
\label{subsec:approx-trisection}

For a general angle $\theta$, exact Euclidean trisection is not available by
straightedge and compass~\cite{wantzel}.  A specific constructible approximation
associated with Baker's construction~\cite{baker-trisection,baker-proportional},
presented interactively by Rostamian~\cite{rostamian}, is
\begin{equation}
  D(\theta)
  =2\arcsin\!\left(\frac43\sin\frac{\theta}{8}\right).
  \label{eq:Dtheta}
\end{equation}
Its local expansion is
\begin{equation}
  D(\theta)
  =\frac{\theta}{3}
   +\frac{7}{10368}\theta^3
   +\Order(\theta^5),
  \label{eq:Dexpansion}
\end{equation}
so the initial angular error is cubic.  The companion trisection work applies
the same chord-transfer principle to the residual and obtains a local cubic
cascade.  Each finite stage is constructible, while the exact trisection value
is reached only in the limiting refinement.

The central angle of the ideal regular $N$-gon is $2\pi/N$.  Under a tripling
schedule, the ideal analytic sequence uses $N_j=N_0 3^j$.  At a finite
constructive stage, let the available half-angle be
\begin{equation}
  \widetilde x_N=\frac{\pi}{N}+\eta_N,
  \label{eq:angle-perturbation}
\end{equation}
where $\eta_N$ is the half-angle residual.  Define the corresponding chord-based
semiperimeter estimator by
\begin{equation}
  \widetilde P_N:=N\sin(\widetilde x_N).
  \label{eq:Ptilde}
\end{equation}
This notation does not assert that the approximate central angles close to an
exact regular polygon; it records the constructible chord length multiplied by
the intended side count.

\begin{lemma}[Transfer of angular residual to perimeter error]
\label{lem:transfer}
For every $N\ge2$,
\begin{equation}
  \abs{\widetilde P_N-P_N}\le N\abs{\eta_N}.
  \label{eq:perturbation-bound}
\end{equation}
Consequently:
\begin{enumerate}
  \item if $\eta_N=o(N^{-3})$, then
        $\widetilde P_N=P_N+o(N^{-2})$, so the raw $N^{-2}$ scale law is
        unchanged;
  \item if $\eta_N=o(N^{-5})$ at every scale entering
        $\Rop_m[\widetilde P]_N$, then the fourth-order cancellation law is
        unchanged.
\end{enumerate}
\end{lemma}

\begin{proof}
By the mean-value theorem and $\abs{\cos t}\le1$,
\[
  \abs{\sin(\pi/N+\eta_N)-\sin(\pi/N)}\le\abs{\eta_N}.
\]
Multiplication by $N$ yields \eqref{eq:perturbation-bound}.  The two stated
conditions follow by comparing $N\abs{\eta_N}$ with $N^{-2}$ and $N^{-4}$,
respectively.  Fixed linear weights in $\Rop_m$ do not alter the asymptotic
order.
\end{proof}

\begin{remark}[Why ``machine precision'' is not a theorem]
A fixed absolute angular error will eventually dominate an increasingly small
polygonal truncation error.  The relevant requirement is therefore not that the
residual be described informally as ``machine precision,'' but that the working
precision and the number of residual refinements be chosen so that
\eqref{eq:perturbation-bound} remains below the target asymptotic error.  In an
arbitrary-precision implementation, this condition can be enforced at every
finite stage.
\end{remark}


\subsection{A certified constructible refinement schedule}
\label{subsec:ladder}

Lemma~\ref{lem:transfer} is conditional.  It states how small the angular
residual must be for each asymptotic law to survive, but it does not establish
that any particular constructible ladder attains that size.  Since the title of
this paper claims a constructive lane, that gap must be closed.  This subsection
closes it.  Throughout we work in an \emph{exact-operation model}: angle and
chord constructions are assumed to be performed without implementation
tolerance, and the indicated small-angle branches of the inverse trigonometric
functions are selected.  Physical drawing tolerance is a separate question and
is not modelled here.

Write $M_i=6\cdot3^i$, let $x_i^\ast=\pi/M_i$ denote the exact half-angle at
level $i$, and let $\widehat x_0=\pi/6$ be the exactly constructible hexagonal
start.  For an input half-angle $x$ put
\begin{equation}
  s(x)=\arcsin\!\left(\tfrac43\sin\tfrac x4\right),
  \qquad
  d(x)=s(x)-\frac x3,
  \qquad
  F(z)=z-s(3z).
  \label{eq:seed-residual}
\end{equation}
Here $s(x)=\tfrac12 D(2x)$ is the half-angle form of the seed
\eqref{eq:Dtheta}, so \eqref{eq:Dexpansion} gives
$d(x)=\tfrac{7}{2592}x^3+\Order(x^5)$, and $F$ is the residual-transfer map
obtained by applying the same seed to three times the measured defect.  The
defect is accessible to the construction without any numerical oracle for
$x/3$: comparing the parent angle $2x$ with three copies of the seed yields the
signed residual $2x-6s(x)=-6d(x)$.  The two maps are related by
\begin{equation}
  F(z)=-d(3z),
  \label{eq:F-identity}
\end{equation}
immediately from \eqref{eq:seed-residual}; this collapses the $F$ half of
Lemma~\ref{lem:seed-bounds} to the $d$ half and is what the formal development
uses.  After $k_i$ refinements at level $i$,
\begin{equation}
  \widehat x_i=\frac{\widehat x_{i-1}}{3}+r_i,
  \qquad
  r_i=F^{\circ k_i}\!\bigl(d(\widehat x_{i-1})\bigr),
  \label{eq:ladder}
\end{equation}
where $k_i=0$ means the seed is used with no refinement.  Set
$e_i=\widehat x_i-x_i^\ast$.

\begin{proposition}[Exact inheritance identity]
\label{prop:inheritance}
Every ladder of the form \eqref{eq:ladder} satisfies
\begin{equation}
  e_i=\frac{e_{i-1}}{3}+r_i,
  \qquad
  M_Le_L=M_0e_0+\sum_{i=1}^{L}M_ir_i.
  \label{eq:inheritance}
\end{equation}
\end{proposition}

\begin{proof}
Since $x_i^\ast=x_{i-1}^\ast/3$, subtracting $x_i^\ast$ from \eqref{eq:ladder}
gives the first identity.  Multiplying by $M_i=3M_{i-1}$ and telescoping gives
the second.
\end{proof}

The content of \eqref{eq:inheritance} is a negative one: the refinement $F$ is
applied to the fresh local defect $d(\widehat x_{i-1})$ and therefore does
\emph{not} act on the inherited term $e_{i-1}/3$.  Any argument that assumes a
refinement stage contracts the \emph{whole} residual is assuming a mechanism
that \eqref{eq:ladder} does not provide, and a recurrence of the form
$e_i\le C(e_{i-1}/3+aM_i^{-3})^3$ is not the recurrence of this ladder.

\begin{lemma}[Elementary seed and cubic bounds]
\label{lem:seed-bounds}
For $0<x\le\tfrac32$,
\begin{equation}
  0<d(x)\le\frac{x^3}{81},
  \label{eq:seed-bounds}
\end{equation}
and consequently, for $\abs z\le\tfrac12$,
\begin{equation}
  \operatorname{sgn}F(z)=-\operatorname{sgn}z\quad(z\ne0),
  \qquad
  \abs{F(z)}\le\frac{\abs z^3}{3}.
  \label{eq:F-bounds}
\end{equation}
\end{lemma}

\begin{proof}
Granted \eqref{eq:seed-bounds}, both parts of \eqref{eq:F-bounds} follow from
the identity \eqref{eq:F-identity}.  At $z=0$ both sides of the magnitude claim
vanish, since $s(0)=0$ and hence $F(0)=0$.  For $0<\abs z\le\tfrac12$ we have
$0<3\abs z\le\tfrac32$, so $F(z)=-d(3z)$ is nonzero of sign opposite to $z$,
and, $d$ being odd in the relevant sense through \eqref{eq:F-identity},
\[
  \abs{F(z)}=d(3\abs z)\le\frac{(3\abs z)^3}{81}=\frac{\abs z^3}{3}.
\]
For \eqref{eq:seed-bounds}, note that $x/3$ and $x/3+x^3/81$ both lie in
$[0,\pi/2]$ for $0<x\le\tfrac32$, and that $\arcsin$ is increasing; since
$s(x)=\arcsin\bigl(\tfrac43\sin(x/4)\bigr)$, the claim is therefore equivalent
to the pair
\begin{equation}
  \sin\frac x3<\frac43\sin\frac x4,
  \qquad
  \frac43\sin\frac x4\le\sin\!\left(\frac x3+\frac{x^3}{81}\right),
  \label{eq:two-sine}
\end{equation}
the first giving $d(x)>0$ and the second $d(x)\le x^3/81$.  Both reduce to
polynomial inequalities by the two-sided cubic bound
$\abs{\sin t-(t-t^3/6)}\le\tfrac{5}{96}t^4$, valid for $\abs t\le1$, applied at
$t=x/3$, $t=x/4$ and $t=x/3+x^3/81$; on $(0,\tfrac32]$ the resulting gaps are
at least $1.33\times10^{-3}x^3$ and $6.24\times10^{-3}x^3$ respectively.
\end{proof}

\begin{remark}[On the range and the constant]
\label{rem:seed-range}
The interval $(0,\tfrac32]$ is wider than the $(0,\pi/6]$ one might expect from
the construction, and deliberately so: since $3\cdot\tfrac12=\tfrac32$, it is
exactly what makes \eqref{eq:F-bounds} available on all of $\abs z\le\tfrac12$
through \eqref{eq:F-identity}, with no separate argument for $F$.  The constant
$1/81$ is not sharp.  The true leading behaviour is
$d(x)=\tfrac{7}{2592}x^3+\Order(x^5)$, and at the right endpoint
\[
  \frac{81\,d(3/2)}{(3/2)^3}=0.2451186\ldots,
\]
numerical sampling on $(0,\tfrac32]$ finding no larger value; we do not prove
that the endpoint is the maximum, and do not need to.  A sharper constant would
shorten the schedules of Remarks~\ref{rem:ladder-instance}
and~\ref{rem:eleven-instance} but is not needed.
\end{remark}

\begin{proposition}[One refinement per level does not suffice]
\label{prop:one-refinement-fails}
Let $\widehat x_0=\pi/6$ and $k_i=1$ for every $i\ge1$.  Then $M_Le_L$ converges
to a strictly negative constant $B$, so that $e_L=\Theta(M_L^{-1})$ and
\[
  M_L\sin(\widehat x_L)\longrightarrow\pi+B\ne\pi.
\]
\end{proposition}

\begin{proof}
By Lemma~\ref{lem:seed-bounds} each $r_i=F(d(\widehat x_{i-1}))$ is strictly
negative, and $\abs{r_i}\le x^9/(3\cdot81^3)$ for $x=\widehat x_{i-1}$, which is
smaller than $x/3$; hence the inputs remain positive and bounded by $\pi/6$ by
induction, and $\widehat x_i<\widehat x_{i-1}/3$ gives the upper range.  Since
$\widehat x_i\le\pi/M_i$ we have $\abs{r_i}=\Order(M_i^{-9})$, so
$\sum_iM_ir_i$ converges absolutely to some $B<0$; with $e_0=0$,
Proposition~\ref{prop:inheritance} gives $M_Le_L\to B$.  Finally
$M_L\sin(\widehat x_L)-\pi=M_L\bigl(\sin\widehat x_L-\widehat x_L\bigr)+M_Le_L$
and the first term is $\Order(M_L^{-2})$.
\end{proof}

A high-precision evaluation of the hexagon-seeded one-refinement ladder gives
\[
  M_Le_L\longrightarrow B=-7.95113571713\times10^{-11},
\]
stable to the digits shown from $L=3$ onward.  This figure is a numerical
diagnostic; the sign and the nonvanishing of the limit are established
analytically by Proposition~\ref{prop:one-refinement-fails}.  The bias is far
smaller than the $7/16$ leading-coefficient shift of the \emph{unrefined}
single-level seed, but it does not vanish, and an $\Theta(M_L^{-1})$ half-angle
residual violates both hypotheses of Lemma~\ref{lem:transfer}.

The repair is to stop asking for a single infinite ladder with a fixed
refinement count and instead fix the target depth first.

\begin{theorem}[Finite-target certified ladder]
\label{thm:finite-target}
Fix a level $L\ge1$ and a desired total scaled error $\tau_L>0$.  Suppose that
for every $i=1,\ldots,L$ the refinement count $k_i=k_i(L)$ in \eqref{eq:ladder}
is chosen so that
\begin{equation}
  \abs{r_i}\le\frac{\tau_L}{LM_i}.
  \label{eq:budget}
\end{equation}
Then the hexagon-seeded ladder satisfies
\begin{equation}
  \abs{\widehat x_L-\frac{\pi}{M_L}}\le\frac{\tau_L}{M_L},
  \qquad
  \abs[\big]{M_L\sin(\widehat x_L)-M_L\sin(\pi/M_L)}\le\tau_L.
  \label{eq:finite-target}
\end{equation}
\end{theorem}

\begin{proof}
Apply \eqref{eq:inheritance} with $e_0=0$ and the triangle inequality:
$M_L\abs{e_L}\le\sum_{i=1}^{L}M_i\abs{r_i}\le\tau_L$.  The second estimate
follows from $\abs{\sin u-\sin v}\le\abs{u-v}$.
\end{proof}

No smallness of $\tau_L$ is used here: the estimate is a consequence of the
budget \eqref{eq:budget} and the inheritance identity alone.  A restriction on
$\tau_L$ enters only with the explicit schedule below, where the bootstrap needs
it; see Remark~\ref{rem:lean-hypothesis}.

\begin{corollary}[The constructible lane inherits the accelerated law]
\label{cor:lane}
Choosing $\tau_L=M_L^{-5}$ in Theorem~\ref{thm:finite-target} gives
$\abs{e_L}\le M_L^{-6}=o(M_L^{-5})$.  Hence the hypothesis of
Lemma~\ref{lem:transfer}\,(2) holds along the constructible tripling sequence,
and the fourth-order cancellation law of \eqref{eq:R3} is available on that lane
with the same leading constant.
\end{corollary}

The refinement counts can be fixed by an entirely rational recurrence, so no
square roots or transcendental evaluations enter the budget check.  For this
construction we additionally require
\begin{equation}
  \tau_L\le10^{-2},
  \label{eq:tau-small}
\end{equation}
which the bootstrap below uses and which Theorem~\ref{thm:finite-target} itself
does not.  By
Lemma~\ref{lem:seed-bounds}, $k$ refinements starting from a defect majorant
$\abs d\le D<\tfrac12$ satisfy $\abs{F^{\circ k}(d)}\le b_k$, where
$b_0=D$ and $b_{j+1}=b_j^3/3$; iterate until $b_{k_i}\le\tau_L/(LM_i)$.

The bootstrap of \eqref{eq:inheritance} keeps
$\widehat x_{i-1}\le(\pi+\tau_L)/M_{i-1}$, so it suffices to fix any constant
$C$ with
\begin{equation}
  \pi+\tau_L\le C
  \qquad\text{and}\qquad
  C\le9,
  \label{eq:bootstrap-constant}
\end{equation}
the second condition giving $C/M_i\le\tfrac32$ at every level because
$M_i\ge6$.  By Lemma~\ref{lem:seed-bounds} the rational seed majorant
\begin{equation}
  D_i=\frac{1}{81}\left(\frac{C}{M_{i-1}}\right)^{\!3}
  \label{eq:seed-majorant}
\end{equation}
is then admissible and well below $\tfrac12$.  Since $b\mapsto b^3/3$ contracts
on $(0,\tfrac12)$, the process terminates for every finite $L$ and every
$\tau_L>0$.

\begin{remark}[Exact-rational instances]
\label{rem:ladder-instance}
Take $L=10$, so $M_L=354294$, and $\tau_L=10^{-85}$.  The counts returned by the
rational recurrence depend on the bootstrap constant of
\eqref{eq:bootstrap-constant}:
\[
  C=\tfrac{16}{5}:\quad(k_1,\ldots,k_{10})=(4,3,3,3,3,2,2,2,2,2),
  \qquad
  C=\tfrac{401}{100}:\quad(4,3,3,3,3,3,2,2,2,2).
\]
In both cases $\sum_{i=1}^{10}M_ib_{i,k_i}\le\tau_L$ as exact fractions, hence
$\abs{e_{10}}\le10^{-85}/354294$.  Both constants are admissible as
mathematics: $\pi+\tau_L\le\tfrac{16}{5}$ already follows for every
$\tau_L\le10^{-2}$ from the classical $\pi<22/7$, since
$22/7+10^{-2}<\tfrac{16}{5}$.  The second schedule is nevertheless the one
carried through the formal development of Section~\ref{sec:lean}, for the
reason given in Remark~\ref{rem:lean-constant}.  The extra refinement at level
$6$ is then necessary and not merely cautious: with $k_6=2$ and
$C=\tfrac{401}{100}$ the budget check fails.  Either schedule certifies the
refinement budget only; neither certifies physical construction tolerances.
\end{remark}

\begin{remark}[Measuring the inherited error directly]
\label{rem:replication}
Proposition~\ref{prop:one-refinement-fails} is a statement about
\eqref{eq:ladder}, in which the correction sees only the local defect.  There is
one further oracle-free mechanism that does expose the inherited error, and we
record it because it delimits exactly what
Proposition~\ref{prop:one-refinement-fails} does and does not rule out.  Because
the exact targets telescope, $3^ix_i^\ast=\pi/6=\widehat x_0$, so comparing
$3^i$ copies of the current half-angle against the constructible hexagonal
half-angle gives
\[
  3^i\widehat x_i-\widehat x_0
  =3^i\bigl(x_i^\ast+e_i\bigr)-\frac{\pi}{6}
  =3^ie_i.
\]
The comparison therefore returns the \emph{total} accumulated error $3^ie_i$,
with no reference to any numerical value of $\pi$, and correcting against it
restores the self-correcting cubic recurrence that \eqref{eq:ladder} lacks.  The
cost is $3^i$ angle replications at level $i$, which is exponential in the
level, so we do not develop this lane further; Theorem~\ref{thm:finite-target}
is the practical route.  We state it to make precise that the obstruction is
one of measurement cost, not of principle.
\end{remark}

% -----------------------------------------------------------
\section{The identity $\pi_g=\pi$}
\label{sec:pig}

Earlier exploratory descriptions treated $\pi_g$ as a numerically observed
geometric constant whose equality with $\pi$ remained open.  In the present
formulation the equality follows directly from limit preservation.

\begin{definition}[Geometric polygonal limit]
For any fixed $N_0\ge3$ and scale factor $m>1$, define
\begin{equation}
  \pi_g:=\lim_{j\to\infty}P_{m^jN_0},
  \label{eq:pig-def}
\end{equation}
provided the limit exists.  One may analogously define accelerated limits from
$\Rop_m[P]_{m^jN_0}$ or from higher N-Series columns.
\end{definition}

\begin{theorem}[Limit preservation]
For the polygonal sequence \eqref{eq:PN},
\begin{equation}
  \pi_g=\pi.
  \label{eq:pig-pi}
\end{equation}
Moreover,
\[
  \lim_{j\to\infty}\Rop_m[P]_{m^jN_0}=\pi,
\]
and every multiscale operator satisfying
\eqref{eq:moment-conditions} has the same limit.
\end{theorem}

\begin{proof}
Since $\sin x/x\to1$ as $x\to0$,
\[
  P_N
  =\pi\frac{\sin(\pi/N)}{\pi/N}
  \longrightarrow\pi.
\]
Hence every subsequence $P_{m^jN_0}$ has limit $\pi$, proving
\eqref{eq:pig-pi}.  The two-scale operator is a fixed linear combination whose
weights sum to one:
\[
  \frac{m^2}{m^2-1}-\frac{1}{m^2-1}=1.
\]
Therefore it preserves any common limit.  The same argument applies to the
multiscale weights because $\sum_jw_j=1$.
\end{proof}

\begin{corollary}[Constructible chord stages]
If $N_j\abs{\eta_{N_j}}\to0$ along a constructible tripling sequence, then
\[
  \widetilde P_{N_j}\longrightarrow\pi.
\]
Thus the limiting value is not a distinct constant; the notation $\pi_g$
records the geometric route by which $\pi$ is approached.
\end{corollary}

\begin{proof}
By \eqref{eq:perturbation-bound},
$\widetilde P_{N_j}-P_{N_j}\to0$, while $P_{N_j}\to\pi$.
\end{proof}

% -----------------------------------------------------------
\section{Numerical verification}

\subsection{Raw doubling and tripling}

Table~\ref{tab:raw-comparison} uses the exact analytic reference
\eqref{eq:PN}.  The entries were evaluated at high working precision; the
printed decimals are rounded only for display.

\begin{table}[htbp]
\centering
\caption{Raw inscribed-polygon errors for matched doubling and tripling stages
starting from $N_0=6$.}
\label{tab:raw-comparison}
\small
\begin{tabular}{@{}r r l r l@{}}
\toprule
Stage & Doubling $N$ & $\pi-P_N$ & Tripling $N$ & $\pi-P_N$ \\
\midrule
0 & 6   & $1.4159265359\times10^{-1}$ & 6    & $1.4159265359\times10^{-1}$ \\
1 & 12  & $3.5764112360\times10^{-2}$ & 18   & $1.5925455585\times10^{-2}$ \\
2 & 24  & $8.9640403086\times10^{-3}$ & 54   & $1.7718924241\times10^{-3}$ \\
3 & 48  & $2.2424505429\times10^{-3}$ & 162  & $1.9690655424\times10^{-4}$ \\
4 & 96  & $5.6070269928\times10^{-4}$ & 486  & $2.1878871714\times10^{-5}$ \\
5 & 192 & $1.4018130433\times10^{-4}$ & 1458 & $2.4309902606\times10^{-6}$ \\
\bottomrule
\end{tabular}
\end{table}

After twelve refinements,
\[
  N_{\mathrm{double}}=6\cdot2^{12}=24576,
  \qquad
  N_{\mathrm{triple}}=6\cdot3^{12}=3188646.
\]
The corresponding errors are
\begin{align*}
  E_{24576}
  &=8.55610234179050\times10^{-9},\\
  E_{3188646}
  &=5.08259791140011\times10^{-13},
\end{align*}
with ratio
\[
  \frac{E_{24576}}{E_{3188646}}
  =16834.1121822748\ldots,
\]
in close agreement with the asymptotic prediction
$(9/4)^{12}=16834.1121960282\ldots$.

\subsection{Fourth-order tripling cancellation}

Table~\ref{tab:R3} evaluates \eqref{eq:R3}.  The ratio of successive errors
approaches $81$, confirming fourth-order behavior along the tripling scale.

\begin{table}[htbp]
\centering
\caption{Errors of the two-scale tripling operator
$\Rop_3[P]_N=(9P_{3N}-P_N)/8$.}
\label{tab:R3}
\small
\begin{tabular}{@{}r r l l@{}}
\toprule
Stage & Base $N$ & $\pi-\Rop_3[P]_N$ & Previous error / current error \\
\midrule
0 & 6    & $2.1705583445\times10^{-4}$ & --- \\
1 & 18   & $2.6970289799\times10^{-6}$ & $80.4796078$ \\
2 & 54   & $3.3320513201\times10^{-8}$ & $80.9420000$ \\
3 & 162  & $4.1139710291\times10^{-10}$ & $80.9935533$ \\
4 & 486  & $5.0790214954\times10^{-12}$ & $80.9992837$ \\
5 & 1458 & $6.2704030692\times10^{-14}$ & $80.9999204$ \\
\bottomrule
\end{tabular}
\end{table}

\subsection{A $\pi$-free chord-ladder verification}

The algebraic recurrence \eqref{eq:cubic-tripling} permits a direct test in
which the approximants are generated without evaluating a trigonometric
function at $\pi/N$.  Representative results from the scale-$3$ computation
are shown in Table~\ref{tab:pi-free}.  The raw column uses the chord ladder
itself; the accelerated column uses the growing N-Series cancellation tableau.

\begin{table}[htbp]
\centering
\caption{Representative $\pi$-free chord-ladder verification.  The reference
value of $\pi$ is used only after the approximants have been generated, to
count matching decimal digits.}
\label{tab:pi-free}
\begin{tabular}{@{}r r r@{}}
\toprule
Tripling levels $L$ & Raw half-perimeter & Accelerated scale-$3$ N-Series \\
\midrule
60  & 57 digits  & 1{,}948 digits \\
95  & 90 digits  & 4{,}705 digits \\
152 & 145 digits & 11{,}199 digits \\
\bottomrule
\end{tabular}
\end{table}

As an implementation cross-check, triangle-seeded and hexagon-seeded versions
of the cubic chord ladder were run independently and their accelerated strings
agreed through $100{,}000$ decimal digits.  In the corresponding certified run,
the largest reported cubic residual was below
$1.3\times10^{-100126}$ and the largest derivative-based root-error bound was
below $4.4\times10^{-100127}$.  These certificates concern the numerical root
solves in \eqref{eq:cubic-tripling}; they do not convert those cubic roots into
straightedge-and-compass constructions.


\subsection{Certified finite-value Richardson enclosure}
\label{subsec:certified-enclosure}

The previous diagonal-difference diagnostic measured numerical stabilization
but did not by itself bound the distance to $\pi$.  The canonical certification
lane of~\cite{baker-certification} replaces that diagnostic as the proof
mechanism.

Two quantities must be distinguished, and we fix notation for them here and use
it throughout: $T_s$ denotes the \emph{exact} tableau value formed from the
analytic polygon inputs, and $\widehat\pi_g$ denotes the \emph{emitted} value of
the finite-precision production computation.  For eleven exact polygon inputs,
let $s=10$ and write
\[
  T_s=\sum_{j=0}^{s}w_jP_{6\cdot3^j},
\]
where the rational weights are generated by the Richardson recurrence.  They
satisfy, exactly,
\[
  \sum_{j=0}^{s}w_j=1,
  \qquad
  \sum_{j=0}^{s}\frac{w_j}{(6\cdot3^j)^{2m}}=0
  \quad(1\le m\le s).
\]

\begin{theorem}[Certified canonical scale-$3$ value]
\label{thm:certified-canonical-value}
Let $\widehat\pi_g$ be the value emitted by the canonical production computation
using $120$ decimal digits and the eleven chord levels $L=0,\ldots,10$.  Exact
rational cubic-root intervals and exact rational affine propagation through the
Richardson tableau give
\[
  \abs{\widehat\pi_g-T_{10}}
  \le
  B_{\rm num}
  =1.60292123843046314\times10^{-89}.
\]
Using $U=22/7$ and the alternating sine-series remainder gives
\[
  \abs{T_{10}-\pi}
  \le
  B_{\rm an}
  =3.40628574228490542\times10^{-81}.
\]
Therefore
\begin{equation}
  \abs{\widehat\pi_g-\pi}
  \le B_{\rm num}+B_{\rm an}
  <3.406285758314118\times10^{-81}
  <10^{-80}.
  \label{eq:certified-r9-bound}
\end{equation}
Thus the canonical finite value is certified to at least $80$ decimal places
of absolute accuracy without consulting a numerical $\pi$ oracle.
\end{theorem}

\begin{proof}
For every $N\ge6$, truncate the sine expansion after the $N^{-20}$ term.  Since
$0<\pi/N<1$, the alternating terms decrease, and
\[
  \abs{r_{10}(N)}
  \le
  \frac{U^{23}}{23!N^{22}}.
\]
The exact moment identities preserve the constant term and cancel the ten modes
\[
  N^{-2},\ N^{-4},\ \ldots,\ N^{-20}.
\]
Hence
\[
  \abs{T_{10}-\pi}
  \le
  \frac{U^{23}}{23!}
  \sum_{j=0}^{10}\frac{\abs{w_j}}{(6\cdot3^j)^{22}}
  =B_{\rm an}.
\]
The weights are the Lagrange basis at the nodes $x_j=(6\cdot3^j)^{-2}$
evaluated at the origin, $w_i=\prod_{k\ne i}x_k/(x_k-x_i)$, so a referee can
regenerate them from this paper alone.  Exact rational evaluation gives
$\sum_jw_j=1$ with moments $1$ through $10$ vanishing identically,
$\max_j\abs{w_j}=1.1408\ldots$, $\sum_j\abs{w_j}=1.2852\ldots$, and
$\sum_{j=0}^{10}\abs{w_j}/(6\cdot3^j)^{22}=3.2085\times10^{-70}$.  The tableau
is therefore \emph{moderately} conditioned rather than amplification-free:
because $\sum_jw_j=1$ while $\sum_j\abs{w_j}>1$, independent input-error bounds
are magnified in the worst case by the factor $1.2852\ldots$.  The absolute
weights are strongly concentrated on the finest two levels,
$\abs{w_{10}}=1.1408\ldots$ and $\abs{w_9}=0.1426\ldots$, with
$\abs{w_j}\le2.45\times10^{-6}$ for every $j\le7$, so that factor is governed
almost entirely by those two levels.  All printed
bounds are rounded upward in the last displayed digit, so each displayed value
is itself a valid bound.
The recursive cubic coefficient is enclosed at each level by exact dyadic
rational endpoints.  Monotonicity of the selected root with respect to the
coefficient permits interval-valued input propagation, and sign-aware exact
rational affine arithmetic encloses the complete Richardson tableau.  The
production value $\widehat\pi_g$ is itself an exact binary rational, so its
distance from that tableau interval is bounded by $B_{\rm num}$.  The triangle inequality
then yields \eqref{eq:certified-r9-bound}.
\end{proof}


\begin{corollary}[Certified constructible-lane tableau value]
\label{cor:eleven-input}
Theorem~\ref{thm:certified-canonical-value} draws its eleven polygon inputs from
the exact algebraic recurrence \eqref{eq:cubic-tripling}, and $B_{\rm num}$
accounts for the root solves only.  Suppose instead that every input is supplied
by the constructible ladder \eqref{eq:ladder}, so that
$\widetilde P_{M_j}=M_j\sin(\widehat x_j)$ with $M_j=6\cdot3^j$, and put
$\widetilde T_{10}=\sum_{j=0}^{10}w_j\widetilde P_{M_j}$.  Fix a construction
allowance $B_{\rm con}>0$ and run a \emph{single} ladder to level $10$ under the
budget of Theorem~\ref{thm:finite-target} with
\begin{equation}
  \tau=\frac{B_{\rm con}}{\sum_{j=0}^{10}\abs{w_j}}.
  \label{eq:eleven-budget}
\end{equation}
Then $M_j\abs{e_j}\le\tau$ holds for all $j\le10$ simultaneously, and
\begin{equation}
  \abs[\big]{\widetilde T_{10}-T_{10}}
  \le\sum_{j=0}^{10}\abs{w_j}\,M_j\abs{e_j}
  \le B_{\rm con},
  \qquad
  \abs[\big]{\widetilde T_{10}-\pi}\le B_{\rm con}+B_{\rm an}.
  \label{eq:eleven-total}
\end{equation}
The two contributions are the construction allowance and the analytic
truncation only.  A finite-precision \emph{evaluation} of $\widetilde T_{10}$
would carry a third term in the role of $B_{\rm num}$, bounding the arithmetic
of the tableau itself; \eqref{eq:eleven-total} is stated for exact evaluation of
the constructed inputs.
\end{corollary}

\begin{proof}
Proposition~\ref{prop:inheritance} with $e_0=0$ gives
$M_je_j=\sum_{i=1}^{j}M_ir_i$ for every $j\le10$, so the single budget
$\sum_{i=1}^{10}M_i\abs{r_i}\le\tau$ supplied by
Theorem~\ref{thm:finite-target} bounds $M_j\abs{e_j}$ by $\tau$ at every level
at once rather than only at the top level.
Lemma~\ref{lem:transfer} gives
$\abs{\widetilde P_{M_j}-P_{M_j}}\le M_j\abs{e_j}$, so the triangle inequality
and \eqref{eq:eleven-budget} yield the first estimate in
\eqref{eq:eleven-total}; the second combines it with
Theorem~\ref{thm:certified-canonical-value}.
\end{proof}

\begin{remark}[Cost of the eleven-input budget]
\label{rem:eleven-instance}
Take $B_{\rm con}=10^{-81}$, which is far inside \eqref{eq:tau-small} once
divided by $\sum_j\abs{w_j}$, so the explicit schedule applies.  The allowance
does not leave the analytic bound unchanged; it is \emph{added} to it, giving
\[
  \abs[\big]{\widetilde T_{10}-\pi}
  \le B_{\rm con}+B_{\rm an}
  =4.40628574228\ldots\times10^{-81}
  <10^{-80},
\]
so the $80$-decimal threshold of
Theorem~\ref{thm:certified-canonical-value} is preserved on the constructible
lane, at the cost of raising the certified bound by rather less than a factor of
two.  This choice gives $\tau=7.7808\ldots\times10^{-82}$.  The rational
recurrence of
Section~\ref{subsec:ladder} then returns the refinement counts
\[
  (k_1,\ldots,k_{10})=(4,3,3,3,3,2,2,2,2,2)
\]
for either admissible constant, the allowance here being loose enough that the
choice of $C$ does not change them; this is \emph{not} the schedule of
Remark~\ref{rem:ladder-instance} at $C=\tfrac{401}{100}$, whose much smaller
$\tau_L=10^{-85}$ forces an extra refinement at level $6$.  The resulting
exact-rational budget does depend on $C$:
\[
  \sum_{i=1}^{10}M_ib_{i,k_i}
  =5.9378\times10^{-88}\ \ \bigl(C=\tfrac{16}{5}\bigr),
  \qquad
  2.6271\times10^{-85}\ \ \bigl(C=\tfrac{401}{100}\bigr),
\]
consuming about $10^{-6}$ and $3.4\times10^{-4}$ of the allowance
respectively.  The margins are this wide because $b\mapsto b^3/3$ triples the
exponent at each refinement, so a single extra refinement at the coarsest
levels buys many orders of magnitude.  The eleven-input constructible
certificate is therefore essentially free once
Theorem~\ref{thm:finite-target} is available.

The formal development of Section~\ref{sec:lean} certifies a \emph{different and
tighter} instance of Corollary~\ref{cor:eleven-input}, namely
$B_{\rm con}=1.3\times10^{-85}$ at $\tau_L=10^{-85}$, which is
\eqref{eq:lean-tableau}; the present instance with $B_{\rm con}=10^{-81}$ is
checked in exact rational arithmetic but is not the one carried through Lean.
As always this certifies the refinement budget in the exact-operation model,
not physical construction tolerance.
\end{remark}

\begin{remark}[Role of the finite-value check]
The reproducibility report~\cite{baker-certification} records this theorem as
the ninth of its numbered checks.  The change in the last tableau diagonal
remains useful as a convergence diagnostic, but it is not used as an error
bound; the cross-domain identity checks provide additional reproducibility
evidence, and the finite-value enclosure proved here carries the rigorous
certification.
\end{remark}

\subsection{Arbitrary-precision digit comparison}

The retained large-scale verification run used the \texttt{mpmath}
arbitrary-precision library~\cite{mpmath} with a working precision set above the
target output length.  Intermediate values were not deliberately truncated to displayed
decimal strings; like every finite-precision computation, however, arithmetic
was rounded according to the selected \texttt{mpmath} context.  The computed
N-Series value was compared digitwise with a reference expansion generated by an
separately generated Chudnovsky-series reference~\cite{chudnovsky}.
The historical test report records agreement in the first
$1{,}000{,}000$ decimal digits. The exact reference implementation, version,
and output hashes have not been recovered for this run, so this comparison
is reported as historical implementation evidence rather than a fully
reproducible independent certification.  Unlike Theorem~\ref{thm:certified-canonical-value}, this
large-target run is not presently accompanied by a million-digit a priori
Richardson enclosure.

This test has a specific interpretation:
\begin{enumerate}
  \item it detects implementation, indexing, coefficient, and precision errors
        that would alter the decimal prefix;
  \item it provides strong reproducibility evidence when accompanied by the
        script, parameter record, software version, output hash, and reference
        hash;
  \item it does not prove $\pi_g=\pi$, because that equality is already supplied
        by the limit-preservation theorem;
  \item it does not by itself establish superiority over Chudnovsky or other
        state-of-the-art algorithms in arithmetic complexity or wall-clock
        time.
\end{enumerate}

A publication archive should therefore retain the exact target precision,
guard-digit setting, N-Series depth, coefficient-generation method, software
version, elapsed time, and cryptographic hashes of the compared outputs.

% -----------------------------------------------------------
\section{Scope, limitations, and open problems}

\subsection{Scale factor versus geometric constructibility}

The analytic raw scale law and the extrapolation operator are valid for any
fixed real scale factor $m>1$ for which the required sequence values are
available.  A finite straightedge-and-compass realization is a separate
question.  In particular, the existence of an analytic $m$-scale formula does
not imply exact constructibility of a general $m$-section of an arbitrary
angle.

\subsection{Local versus global cubic refinement}

Equation~\eqref{eq:Dexpansion} is a local statement.  The residual refinement
has cubic convergence after the residual enters a sufficiently small
admissible neighborhood.  A global convergence theorem would additionally
require an invariant angular domain, branch control for inverse trigonometric
operations, and proof that the chosen seed enters the local basin.  No global
claim is required for the finite high-precision computations reported here,
where the residual can be evaluated and bounded directly.

\subsection{Geometric asymptotics beyond the circle}

For the circle, the even-power expansion follows immediately from the sine
series.  Extending the same N-Series mechanism to a broader class of smooth
curves requires a separate theorem establishing the appropriate chordal
perimeter expansion and its regularity assumptions.  The cancellation operator
is conditional on that expansion; it does not create the expansion by itself.

\subsection{Constructed inputs versus algebraic inputs}

Theorem~\ref{thm:certified-canonical-value} is a statement about the algebraic
lane: its eleven inputs come from \eqref{eq:cubic-tripling} and $B_{\rm num}$
bounds the root-solve error, so no construction error is present to omit.
Corollary~\ref{cor:eleven-input} is the corresponding statement for the
constructible lane, and it is a distinct result with its own allowance
$B_{\rm con}$.  A published enclosure should state which lane it refers to; the
two are not interchangeable, and only the latter carries a straightedge-and-compass
interpretation.

\subsection{Computational complexity}

Tripling has a provable per-level advantage over doubling because it advances
more rapidly through the scale variable.  This does not determine total runtime.
A fair computational comparison must use matched implementations, identical
precision targets, equivalent certification requirements, and a cost model for
large-number arithmetic and coefficient generation.  The approximately one-day
million-digit run is therefore reported as a verification computation, not as
a benchmark proving a fixed runtime ratio.  The exact-rational R9 machinery
certifies the canonical $80$-decimal result; extending the same enclosure to a
million-digit target is a separate scale-up problem rather than an unresolved
logical gap in the finite certificate.

% -----------------------------------------------------------
\section{Formal verification in Lean 4}
\label{sec:lean}

A companion formal development verifies key mathematical components of the
N-Series framework in Lean~4~\cite{lean4} with Mathlib~\cite{mathlib}. The
integrated development comprises $28$ modules, organized in dependency groups
that the source labels $A$ through $E$, and
built successfully with Lean~4.36.0-rc2 against Mathlib revision
\texttt{9d581dfff2e41de989072ea63abf3daf69311520}; its release audit
reported 112 named declarations and only the standard foundational axioms
\texttt{propext}, \texttt{Classical.choice}, and \texttt{Quot.sound}.
The precise, publicly available source revision is
\url{https://github.com/CWBaker360/nseries-lean-verification/tree/bf262aaf52a12001ca64545757d758b507b09a18}.
The successful local build used pinned dependency build caches; an independent
fresh-clone reproduction on a separate machine is not claimed.

Table~\ref{tab:lean} summarizes representative results of the companion
formalization. The theorem and corollary identifiers in the original standalone
verification note follow that note's own numbering; this table identifies the
corresponding mathematical content without imposing that numbering on the
present manuscript.

The companion development uses its own notation, which maps onto the notation of
this paper as follows.  Its scale parameter $k$ plays the role of the scale
factor $m$ of \eqref{eq:Rop}; its $\widehat P_N^{(k,3)}$ is the three-scale
member of the multiscale family $\Rmulti_{m,s}$ of \eqref{eq:multi-operator}
with $m=k$ and $s=2$; and its $\pi_g^{(k)}$ is the accelerated limit
$\lim_j\Rop_m[P]_{m^jN_0}$ of Section~\ref{sec:pig}, written there as $\pi_g$
when the scale is fixed.  ``Intrinsic'' refers to the $\pi$-free polygonal
definition displayed below, and ``dyadic'' to the specialization $k=2$.

\begin{table}[htbp]
\centering
\small
\begin{tabular}{@{}p{0.42\linewidth}p{0.51\linewidth}@{}}
\toprule
Verified mathematical content & Lean declaration \\
\midrule
Sine expansion and explicit remainder & \texttt{nseries\_sine\_remainder\_bound} \\
General two-scale bound & \texttt{two\_scale\_general\_bound} \\
Polygonal Richardson bound & \texttt{richardson\_error\_bound} \\
Accelerated limit equals $\pi$ & \texttt{acceleratedPi\_eq\_pi} \\
Three-scale bound and sharp term & \texttt{three\_scale\_error\_bound}, \texttt{three\_scale\_sharp} \\
Sharp dyadic three-scale term & \texttt{dyadic\_three\_scale\_sharp} \\
General moment cancellation & \texttt{general\_moment\_cancellation} \\
Moment-weight existence and uniqueness & \texttt{nseries\_weights\_existsUnique}, \texttt{moment\_weights\_unique} \\
Lagrange weights & \texttt{lagrange\_weights\_moments}, \texttt{lagrange\_weight\_formula} \\
Trisection cubic root uniqueness & \texttt{chordHalf\_trisection\_unique} \\
Intrinsic polygonal limit & \texttt{polygonalPi\_eq\_pi} \\
Intrinsic two-scale bound & \texttt{intrinsic\_two\_scale\_bound} \\
Intrinsic dyadic sharp term & \texttt{intrinsic\_dyadic\_three\_scale\_sharp} \\
\midrule
\multicolumn{2}{@{}l}{\emph{Constructible ladder, Section~\ref{subsec:ladder}}}\\
Inheritance identity (Prop.~\ref{prop:inheritance}) & \texttt{err\_succ}, \texttt{ladderM\_mul\_err} \\
Exact hexagonal start & \texttt{ladderM\_mul\_err\_of\_exact\_start}, \texttt{err\_zero\_of\_hex\_start} \\
Per-level budget at every level & \texttt{abs\_ladderM\_mul\_err\_le} \\
Perimeter transfer (Lemma~\ref{lem:transfer}) & \texttt{sin\_ladder\_dist\_le} \\
Finite-target bounds (Thm.~\ref{thm:finite-target}) & \texttt{finite\_target\_ladder} \\
Weighted tableau perturbation & \texttt{constructible\_tableau\_from\_ladder} \\
Eleven-input instance (Cor.~\ref{cor:eleven-input}) & \texttt{constructible\_tableau\_eleven\_inputs} \\
Canonical weights are the Lagrange weights & \texttt{canonicalWeight\_eq\_lagrangeWeight} \\
$\sum_j\abs{w_j}\le13/10$, $\sum_jw_j=1$ & \texttt{canonicalWeight\_abs\_sum\_le}, \texttt{canonicalWeight\_sum} \\
Canonical constructible tableau & \texttt{canonical\_constructible\_tableau} \\
\midrule
\multicolumn{2}{@{}l}{\emph{Seed bounds and the refinement schedule}}\\
Lemma~\ref{lem:seed-bounds}, defect half & \texttt{seedDefect\_pos}, \texttt{seedDefect\_le} \\
The identity \eqref{eq:F-identity} & \texttt{F\_eq\_neg\_seedDefect} \\
Lemma~\ref{lem:seed-bounds}, $F$ half & \texttt{abs\_seedF\_le} \\
Majorant recurrence $b_{j+1}=b_j^3/3$ & \texttt{bseq\_le\_half}, \texttt{cubicIter\_bound} \\
Iteration bound at the seed & \texttt{residual\_bound\_one\_level} \\
Bootstrap of the level input & \texttt{bootstrap\_input} \\
Thm.~\ref{thm:finite-target}, unconditional & \texttt{sched\_certified} \\
Schedule of Rem.~\ref{rem:ladder-instance} & \texttt{canonicalSchedule\_clears\_budget} \\
$M_{10}\abs{e_{10}}\le10^{-85}$ & \texttt{canonical\_ladder\_certified} \\
Per-level residuals & \texttt{canonical\_residual\_bounds} \\
$\abs{\widetilde T_{10}-T_{10}}\le1.3\times10^{-85}$ & \texttt{canonical\_tableau\_certified} \\
Combined, given $B_{\rm an}$ & \texttt{canonical\_tableau\_vs\_pi} \\
\bottomrule
\end{tabular}
\caption{Selected mathematical statements and their Lean formalizations.}
\label{tab:lean}
\end{table}

Four points of scope should be stated precisely.

\emph{Explicit bounds.} The asymptotic statements are verified in
explicit-constant form. For example, the sine expansion is proved as
\[
\left|N\sin\frac{\pi}{N}
-\Bigl(\pi-\frac{\pi^3}{6N^2}+\frac{\pi^5}{120N^4}-\frac{\pi^7}{5040N^6}\Bigr)\right|
\le \frac{\pi^9}{362880\,N^8}
\qquad (N>0),
\]
and the three-scale operator is shown to satisfy
\[
\widehat P_N^{(k,3)}=\pi-\frac{\pi^7}{5040\,k^6N^6}+O(N^{-8})
\]
with an explicit constant. The general two-scale and moment-cancellation results are formalized under an
explicit remainder hypothesis $|q(M)-Q-\cdots|\le B/M^{p}$ rather than a
big-$O$ hypothesis; the Lean statements are therefore quantitative versions
of the corresponding mathematical statements, not verbatim transcriptions of this manuscript.

\emph{The polygonal limit.} The half-perimeter of the regular $n$-gon is
defined in Lean without reference to $\pi$, as
\[
\frac{n}{2}\min_{\substack{z^n=1\\ z\ne 1}}|1-z| \qquad (n\ge 3),
\]
and the limit of this sequence is proved to equal $\pi$. The definition is
$\pi$-free; the proof of the identification uses the classical constant.
This is the raw polygonal limit, distinct from the accelerated limit
$\pi_g^{(k)}$ of the accelerated N-Series construction (the $\pi_g$ of
Section~\ref{sec:pig}), which is verified separately.
The companion development proves acceleration bounds for the analytic sequence
$N\sin(\pi/N)$, and the intrinsic half-perimeter is proved equal to that
sequence for $n\ge 3$. Two of the bounds are also stated and verified
directly for the intrinsic sequence: the two-scale bound
for integer $k\ge 2$, and the sharp bound for the dyadic three-scale approximation. The remaining bounds (the general-$k$ three-scale bound and
the accelerated-limit theorem) have not been restated for the
intrinsic sequence.

\emph{What is not verified.} The digit tables in the numerical-verification section, the high-precision implementation comparisons, and the observed
growth of correct digits with the number of levels are
computational evidence, not formally verified theorems. For the trisection
recurrence, Lean verifies the cubic identity and the uniqueness of its root
in $[0,\tfrac12]$ for $N\ge 2$, not the numerical runs.  That enclosure is
deliberately weaker than
Proposition~\ref{prop:unique-cubic-branch}, which localizes the selected branch
to $(0,c/2)$; the formalized statement is the coarser of the two, and the
sharper interval is not currently part of the Lean development.

\emph{The constructible ladder.}  Section~\ref{subsec:ladder} and
Corollary~\ref{cor:eleven-input} are formalized in a separate standalone
development, audited declaration by declaration to the same three foundational
axioms.  The ledger is as follows.

The residuals are not abstract: the chain is unconditional for the construction
actually described in Section~\ref{subsec:ladder}.

\emph{Verified.}  The inheritance identity \eqref{eq:inheritance} in both forms,
and the vanishing of the inherited term at an exact hexagonal start.  The
perimeter-transfer inequality $\abs{\widetilde P_N-P_N}\le N\abs{\eta_N}$ of
Lemma~\ref{lem:transfer}.  Theorem~\ref{thm:finite-target}, in a form parametric
in the bootstrap constant $C$ of \eqref{eq:bootstrap-constant}, and including
the stronger conclusion that the budget bounds the scaled error at \emph{every}
level $j\le L$ rather than only at the target level --- which is what allows a
single ladder to serve all eleven tableau inputs.  Lemma~\ref{lem:seed-bounds}
in full on the range $(0,\tfrac32]$, together with the identity
\eqref{eq:F-identity} from which its $F$ half is obtained.  The majorant
recurrence $b_{j+1}=b_j^3/3$ and the iteration bound
$\abs{F^{\circ k}(d)}\le b_k$.  The bootstrap: a scaled-error bound of $\tau_L$
places the level input inside $(0,C/M_i]$, so the rational majorant
\eqref{eq:seed-majorant} is admissible.  The schedule of
Remark~\ref{rem:ladder-instance} at $C=\tfrac{401}{100}$, discharged by exact
rational arithmetic at all ten levels, giving
\begin{equation}
  M_{10}\abs[\big]{\widehat x_{10}-\pi/M_{10}}\le10^{-85}
  \label{eq:lean-ladder}
\end{equation}
with no assumed residual budget.  The per-level bounds
$\abs{r_j}\le10^{-85}/(10M_j)$ for $1\le j\le10$, which do \emph{not} follow
from \eqref{eq:lean-ladder} --- a bound on the accumulated error says nothing
about the individual residuals --- and are proved separately from the schedule's
majorants.  The canonical eleven-weight vector, proved equal to the Lagrange
weights $\prod_{k\ne i}x_k/(x_k-x_i)$ at the nodes $x_j=(6\cdot3^j)^{-2}$,
normalized to $\sum_jw_j=1$ and majorized by $\sum_j\abs{w_j}\le13/10$ in exact
rational arithmetic.  And the construction-perturbation component of
Corollary~\ref{cor:eleven-input} at that schedule,
\begin{equation}
  \abs[\big]{\widetilde T_{10}-T_{10}}\le1.3\times10^{-85}.
  \label{eq:lean-tableau}
\end{equation}

\emph{Not formalized.}  Proposition~\ref{prop:one-refinement-fails}, whose
ingredients are now available but which has not been assembled; and
Theorem~\ref{thm:certified-canonical-value}.  The latter is the single remaining
obligation, and the development states it as an explicit hypothesis rather than
leaving it implicit: supplying any $B_{\rm an}$ with
$\abs{T_{10}-\pi}\le B_{\rm an}$ yields
$\abs{\widetilde T_{10}-\pi}\le1.3\times10^{-85}+B_{\rm an}$, and at the value
$B_{\rm an}=3.40628574228490542\times10^{-81}$ of
Theorem~\ref{thm:certified-canonical-value} that total is below $10^{-80}$, so
the $80$-decimal threshold is preserved on the constructible lane.
Theorem~\ref{thm:certified-canonical-value} is \emph{proved} in this paper by
exact rational interval propagation and a closed sine remainder; the absence of
a formalization is not a gap in that proof.

\emph{Weights.} The companion Vandermonde formalization establishes existence and
uniqueness of the moment weights at distinct nodes. Separately, the explicit
Lagrange-basis weights
\[
w_i=\prod_{k\ne i}\frac{-x_k}{x_i-x_k}
\]
are proved to satisfy the moment system. These results certify each finite
moment-cancellation system, not a uniform digit-growth rate as the number of
levels increases.

\begin{remark}[A hypothesis removed by formalization]
\label{rem:lean-hypothesis}
An earlier draft stated Theorem~\ref{thm:finite-target} under
$0<\tau_L\le10^{-2}$.  The Lean development establishes the corresponding bound
with no upper constraint on $\tau_L$ whatever, which exposed the condition as
unused in that theorem's own proof: it is required only by the bootstrap of the
explicit correction-count rule, and it is now stated there as
\eqref{eq:tau-small}.  Corollary~\ref{cor:eleven-input} is correspondingly free
of any restriction on $B_{\rm con}$.  The formal statement also carries the
hypothesis $L\ge1$ of Theorem~\ref{thm:finite-target} explicitly, the budget
\eqref{eq:budget} being undefined at $L=0$.
\end{remark}

\begin{remark}[Why the bootstrap constant is a parameter]
\label{rem:lean-constant}
The constant $C=\tfrac{16}{5}$ is perfectly admissible as mathematics: it
follows from $\pi<22/7$, which is classical and elementary.  The reason the
formal instance uses $C=\tfrac{401}{100}$ is narrower than that.  Among numeric
bounds on $\pi$, the Mathlib revision used here offers only $3<\pi$ and
$\pi<4$ as named results; a formal proof of $\pi<22/7$ would have to be
supplied, and was judged not worth the effort for a constant that only affects
one refinement count.  Hard-coding $\tfrac{16}{5}$ in
\eqref{eq:bootstrap-constant} would meanwhile have made a published constant ---
and with it the schedule of Remark~\ref{rem:ladder-instance} --- depend on which
bounds a library happens to expose.  Theorem~\ref{thm:finite-target} is
accordingly formalized for every $C$ satisfying \eqref{eq:bootstrap-constant},
with the instance supplying $C=\tfrac{401}{100}$ from $\pi<4$ alone; given a
formal $\pi<22/7$ the same development yields the shorter schedule with no
change to any theorem.  This is the second hypothesis the formalization has
relocated rather than merely confirmed; see
Remark~\ref{rem:lean-hypothesis}.
\end{remark}

The reported endpoint axiom audits list only \texttt{propext}, \texttt{Classical.choice},
and \texttt{Quot.sound}; no \texttt{sorry}, \texttt{admit}, or additional axioms were reported
in the audited modules.


% -----------------------------------------------------------
\section{Conclusion}

For the inscribed polygonal approximation
$P_N=N\sin(\pi/N)$, raw tripling and raw doubling inherit their convergence
factors directly from the leading $N^{-2}$ error law.  Tripling reduces the
leading error by a factor approaching $9$ per scale level, compared with $4$
for doubling, and gains $\log 3/\log 2\approx1.58496$ times as many decimal
digits per level.  After twelve matched stages from a hexagon, the exact
tripling reference sequence is approximately $1.6834\times10^4$ more accurate
than the doubling sequence.

The higher-order result comes from scale cancellation.  The operator
\[
  \Rop_3[P]_N=\frac{9P_{3N}-P_N}{8}
\]
annihilates the $N^{-2}$ mode and produces fourth-order convergence.  Higher
N-Series columns are standard multiscale extrapolation combinations applied to
the even-power polygonal expansion.

The geometric interpretation requires care.  Exact analytic regular polygons
serve as the reference sequence; finite Baker stages provide constructible
chord approximations to the required central angles.  The perturbation bound
\[
  \abs{\widetilde P_N-P_N}\le N\abs{\eta_N}
\]
specifies exactly how the angular residual must be controlled for the raw or
accelerated asymptotic law to remain visible.  Under vanishing residual, the
constructible chord sequence approaches the same limit as the Archimedean
sequence.

Accordingly, the geometric constant $\pi_g$ is not numerically distinct from
$\pi$: equality follows from the classical polygonal limit and preservation of
that limit by the N-Series operators.  In addition, the canonical eleven-input
computation is no longer supported only by a last-diagonal stabilization
indicator.  Exact rational root intervals, exact tableau propagation, and a
closed sine-remainder bound prove
\[
  \abs{\widehat\pi_g-\pi}
  <3.406285758314118\times10^{-81}.
\]
The million-digit comparison is retained as a large-scale verification of the
implementation and coefficient pipeline, not as the logical basis for
$\pi_g=\pi$ or as a million-digit formal certificate.

% -----------------------------------------------------------
\section*{Data and code availability}

The manuscript is distributed with the Python source, coefficient files,
reproducibility notes, numerical tables, and checksums in the public research
repository
\begin{center}
\url{https://github.com/CWBaker360/N-Series-pi-acceleration}.
\end{center}
The frozen certification snapshot for Theorem~\ref{thm:certified-canonical-value}
is stored under
\begin{center}
\texttt{verification/certified\_richardson\_v3\_0/}
\end{center}
and contains the oracle-free construction kernel, exact rational certification
module, verification driver, unit tests, canonical outputs, proof report, and
SHA-256 package manifest.  The Lean~4 development is archived separately
at~\cite{nseries-lean}, with the additional frozen November 15 constructive-ladder
sources and audit evidence in the research repository under
\texttt{verification/constructible-ladder-2026-11-15/}. The July 2026
Richardson verification report is also available in the certification directory.
These public GitHub revisions are not represented here as new Zenodo deposits
or DOI-issued manuscript releases.

The original \texttt{SHA256SUMS.txt} at repository root predates the present
verification directories and is not used as a checksum certificate for the
November manuscript. The independently verified numerical package has its own
\texttt{PACKAGE\_MANIFEST.json}, which records one missing supplementary file
(\texttt{REQUIREMENTS.txt}), and the Lean checkpoint has a separate freeze
manifest. Reproducing the numerical result requires the pinned environment and
the verification instructions in those directories.

\section*{Acknowledgments}

The author thanks Rouben Rostamian for his independent analysis of the initial
Baker approximate-trisection construction and for hosting an interactive
geometric presentation of the construction on his UMBC webpage~\cite{rostamian}.  This
acknowledgment does not imply institutional affiliation with UMBC.

% -----------------------------------------------------------
\begin{thebibliography}{99}

\bibitem{archimedes}
Archimedes,
\textit{Measurement of a Circle},
in T. L. Heath (ed.), \textit{The Works of Archimedes},
Cambridge University Press, 1897.

\bibitem{baker-trisection}
W. Baker,
``A Constructible Cubic Error Cascade for Approximate Angle Trisection,''
preprint and research source archive, 2026.

\bibitem{baker-proportional}
W. Baker,
``A Local Cubic Refinement Law for Proportional-Subtended Angle Division,''
preprint and research source archive, 2026.

\bibitem{chudnovsky}
D. V. Chudnovsky and G. V. Chudnovsky,
``The Computation of Classical Constants,''
\textit{Proceedings of the National Academy of Sciences of the United States
of America}, 86(21) (1989), 8178--8182.
\url{https://doi.org/10.1073/pnas.86.21.8178}.

\bibitem{richardson-gaunt}
L. F. Richardson and J. A. Gaunt,
``The Deferred Approach to the Limit. Part I. Single Lattice. Part II.
Interpenetrating Lattices,''
\textit{Philosophical Transactions of the Royal Society of London. Series A},
226 (1927), 299--361.
\url{https://doi.org/10.1098/rsta.1927.0008}.

\bibitem{romberg}
W. Romberg,
``Vereinfachte numerische Integration,''
\textit{Det Kongelige Norske Videnskabers Selskabs Forhandlinger},
28 (1955), 30--36.

\bibitem{rostamian}
R. Rostamian,
``An Angle Trisection,''
University of Maryland, Baltimore County, interactive geometric note.
\url{https://userpages.umbc.edu/~rostamia/Geometry/trisect-baker.html}
(accessed July 2026).

\bibitem{baker-certification}
W. Baker,
``Certified Oracle-Free Construction and Cross-Domain Verification of
Baker's Geometric $\pi_g$ Approximation,''
Version 3.0 reproducibility report, 2026.

\bibitem{lean4}
L. de Moura and S. Ullrich,
``The Lean 4 Theorem Prover and Programming Language,''
in \textit{Automated Deduction -- CADE 28},
Lecture Notes in Computer Science 12699, Springer, 2021, 625--635.
\url{https://doi.org/10.1007/978-3-030-79876-5_37}.

\bibitem{mathlib}
The mathlib Community,
``The Lean Mathematical Library,''
in \textit{Proceedings of the 9th ACM SIGPLAN International Conference on
Certified Programs and Proofs (CPP 2020)},
ACM, 2020, 367--381.
\url{https://doi.org/10.1145/3372885.3373824}.

\bibitem{mpmath}
F. Johansson and others,
\textit{mpmath: a Python library for arbitrary-precision floating-point
arithmetic}.
\url{https://mpmath.org}
(version 1.3.0 for the independently reproduced 120-decimal canonical
certification; version of the earlier million-digit run not recorded).

\bibitem{nseries-lean}
W. Baker,
``N-Series Lean 4 Verification,''
source archive, 2026.
\url{https://github.com/CWBaker360/nseries-lean-verification}.

\bibitem{wantzel}
P. Wantzel,
``Recherches sur les moyens de reconna\^\i tre si un probl\`eme de g\'eom\'etrie
peut se r\'esoudre avec la r\`egle et le compas,''
\textit{Journal de Math\'ematiques Pures et Appliqu\'ees}, 2 (1837), 366--372.

\end{thebibliography}

\end{document}
